% Esquisse d'un programme (cote 119, copy 2): archive cote 119, site page 80 (PDF page 81), typescript page 2.
% typescript. Draft compared: esquisse/pages/esq_119_004.tex (copy 1, site 4).
% AI-transcribed by Claude Opus 5.5 (claude-opus-5-5, adaptive thinking, effort high), Message Batches API, batch esq119_full (msgbatch_015xqDHkTumDysRF1qngZhRJ), 2026-11-13; not yet checked by eye.
\scanpage{119/80}
\mspage{2}

image fausse, ou simplement inadéqua\deleted{n}te, se trouve enfin débusquée
et mise à jour, et que les choses qui semblaient de guingois se
mettent en place, dans l'harmonie mutuelle qui leur est propre.
\qquad Quoi qu'il en soit, l'esquisse qui suit de quelques
thèmes de réflexion des dernières dix ou douze années, tiendra
lieu en même temps d'esquisse de programme de travail pour les
années qui viennent, que je compte consacrer au développement de
ces thèmes, ou au moins de certains d'entre eux. Elle est desti-%
née, d'une part aux collègues du Comité National appelés à statuer
sur ma demande, d'autre part à quelques autres collègues, anciens
élèves, amis, dans l'éventualité où certaines des idées esquissées
ici pourraient intéresser l'un d'entre eux.

2. \qquad Les exigences d'un enseignement universitaire, s'adressant
donc à des étudiants (y compris les étudiants dits ``avancés'') au ba-%
gage mathématique modeste (et souvent moins que modeste), m'ont amené
à renouveler de façon draconienne les thèmes de réflexion à proposer
à mes élèves, et de fil en aiguille et de plus en plus, à moi-même
également. Il m'avait semblé important de partir d'un bagage intuitif
commun, indépendant de tout langage technique censé l'exprimer, bien
antérieur à tout tel langage - il s'est avéré que l'intuition géomé-%
trique et topologique des formes, et plus particulièrement des formes
bidimensionelles, était un tel terrain commun. Il s'agit donc de thèmes
qu'on peut grouper sous l'appellation de ``topologie des surfaces'' ou
``géométrie des surfaces'', étant entendu dans cette dernière appella-%
tion que l'accent principal se trouve sur les propriétés topologiques
des surfaces, ou sur les aspects combinatoires qui en constituent
l'expression technique la plus terre-à-terre, et non sur les aspects
différentiels, voire conformes, riemaniens, holomorphes et (de là)
l'aspect ``courbes algébriques complexes''. Une fois ce dernier pas
franchi cependant, voici soudain la géométrie algébrique (mes anciennes
amours !) qui fait irruption à nouveau, et ce par les objets qu'on
peut considérer comme les pierres de construction ultimes de toutes
les autres variétés algébriques. Alors que dans mes recherches d'avant

\margin{80}
